\chapter{Implémentation et Tests}

\section{Implémentation matérielle}
\subsection{Assemblage mécanique}
\todo{Décrire l'assemblage : montage moteur, table d'indexation, profilés aluminium, supports 3D. Photos du système assemblé.}

\placeholderfigure{Système assemblé VisionGuard}{Système VisionGuard assemblé — vue d'ensemble}

\subsection{Câblage électronique}
\todo{Décrire le câblage complet : PSU, TMC2208, STM32, RPi5, LEDs éclairage. Schéma de câblage.}

\placeholderfigure{Schéma câblage}{Schéma de câblage électronique du système VisionGuard}

\section{Implémentation logicielle}
\subsection{Environnement de développement}
\todo{Ubuntu 22.04, VS Code, pyenv Python 3.11.9, STM32CubeIDE 2.1.1, dépendances Python.}

\subsection{Firmware STM32}
\todo{Décrire l'implémentation firmware : configuration CubeMX, timer STEP, UART HAL polling, tests unitaires réalisés.}

\subsection{Scripts Python}
\todo{Décrire l'implémentation des scripts : rpi5\_server.py, collect\_orchestrator.py, structure du dépôt GitHub.}

\section{Tests et validation}
\subsection{Tests du sous-système mécanique}
\todo{Tests d'indexation : précision angulaire mesurée, répétabilité, vitesse. Résultats numériques.}

\subsection{Tests du sous-système vision}
\todo{Tests optiques : résolution effective sur pièce M00, détection visuelle de défauts $\geq 0.8$~mm, conditions d'éclairage.}

\subsection{Tests d'intégration système}
\todo{Tests de la chaîne complète : motion + capture + transfert. Mesure du temps de cycle par veine.}

\subsection{Entraînement et validation du modèle IA}
\todo{Décrire le dataset collecté (volume, répartition classes), l'entraînement, les métriques obtenues (précision, rappel, F1-score, matrice de confusion).}

\section{Conclusion}
\todo{Synthèse des résultats de tests et transition vers le chapitre résultats/validation.}
